\section*{Цель работы}
Практическое ознакомление с синусоидальными режимами в простых $RL$-, $RC$- и $RLC$-цепях, анализ амплитудных и фазовых соотношений между токами и напряжениями.

\section*{Краткие теоретические сведения}
В установившемся синусоидальном режиме комплексные сопротивления индуктивности и емкости зависят от частоты:
\begin{equation*}
    Z_L = j\omega L = \omega L e^{j90^\circ}; \quad Z_C = \frac{1}{j\omega C} = \frac{1}{\omega C} e^{-j90^\circ}.
\end{equation*}
Для последовательной $RLC$-цепи полное комплексное сопротивление определяется выражением:
\begin{equation*}
    Z = R + Z_L + Z_C = R + j\left( \omega L - \frac{1}{\omega C} \right).
\end{equation*}
Модуль полного сопротивления и фазовый сдвиг $\varphi$ (аргумент $Z$) рассчитываются как:
\begin{equation*}
    |Z| = \sqrt{R^2 + \left( \omega L - \frac{1}{\omega C} \right)^2}; \quad \varphi = \text{arctg} \frac{\omega L - 1/(\omega C)}{R}.
\end{equation*}
Действующее значение тока в цепи определяется по закону Ома:
\begin{equation*}
    I = \frac{U_0}{\sqrt{R^2 + \left( \omega L - 1/(\omega C) \right)^2}}.
\end{equation*}
При условии $\omega L = 1/(\omega C)$ реактивная составляющая сопротивления обращается в ноль, и в цепи наступает резонанс напряжений. Ток и напряжение при этом совпадают по фазе ($\varphi = 0$).

\section*{Экспериментальные исследования}

Для выполнения работы используются: генератор сигналов (ГС), осциллограф, модуль с пассивными элементами ($R, L, C$).

\subsection*{1. Исследование $RC$-цепи}

\begin{figure}[H]
    \centering
    \includegraphics[width=0.6\linewidth]{figs/scheme_rc.pdf}
    \caption{Схема для исследования $RC$-цепи}
    \label{fig:scheme_rc}
\end{figure}

\begin{enumerate}
    \item Соберите схему, представленную на \cref{fig:scheme_rc}. Резистор $R_{01} = 50$ Ом включите параллельно выходу генератора для согласования.
    \item Установите на выходе ГС напряжение $U_0 = 2$ В и частоту $f = 7.5$ кГц.
    \item Измерьте мультиметрами действующие значения тока $I$ и напряжений $U_R$, $U_C$.
    \item С помощью осциллографа снимите осциллограмму входного напряжения и тока (напряжения на резисторе). Определите фазовый сдвиг $\varphi_{\text{осц}}$ между ними.
        \begin{itemize}
            \item \textit{Примечание:} Ток в емкости опережает напряжение, поэтому осциллограмма тока должна быть сдвинута влево относительно напряжения.
        \end{itemize}
    \item Повторите измерения (пп. 2--4) при частоте $f = 15$ кГц.
    \item Результаты занесите в Таблицу 1 файла протокола.
\end{enumerate}

\subsection*{2. Исследование $RL$-цепи}

\begin{figure}[H]
    \centering
    \includegraphics[width=0.6\linewidth]{figs/scheme_rl.pdf}
    \caption{Схема для исследования $RL$-цепи}
    \label{fig:scheme_rl}
\end{figure}

\begin{enumerate}
    \item Соберите схему, представленную на \cref{fig:scheme_rl}.
    \item Установите напряжение $U_0 = 2$ В.
    \item Проведите измерения тока $I$, напряжений $U_R, U_L$ и фазового сдвига $\varphi_{\text{осц}}$ для двух частот: $f = 7.5$ кГц и $f = 15$ кГц.
    \item Результаты занесите в Таблицу 1 файла протокола.
\end{enumerate}

\subsection*{3. Исследование $RLC$-цепи (Резонанс напряжений)}

\begin{figure}[H]
    \centering
    \includegraphics[width=0.6\linewidth]{figs/scheme_rlc.pdf}
    \caption{Схема для исследования $RLC$-цепи}
    \label{fig:scheme_rlc}
\end{figure}

\begin{enumerate}
    \item Соберите схему, представленную на \cref{fig:scheme_rlc}.
    \item Установите на выходе ГС напряжение $U_0 = 2$ В.
    \item Переключите осциллограф в режим «X-Y» (фигуры Лиссажу).
    \item Изменяя частоту ГС в пределах $3 \dots 12$ кГц, добейтесь резонанса.
        \begin{itemize}
            \item Критерий резонанса: эллипс на экране вырождается в прямую линию (фазовый сдвиг $\varphi = 0$).
        \end{itemize}
    \item Определите резонансную частоту $f_0 = 1/T_0$ по периоду синусоиды (переключив осциллограф обратно в режим развертки по времени).
    \item Выполните измерения $I, U_R, U_C, U_L$ и $\varphi_{\text{осц}}$ для трех частот:
        \begin{itemize}
            \item $f = f_0$ (резонанс);
            \item $f = 2f_0$;
            \item $f = 0.5f_0$.
        \end{itemize}
    \item Результаты занесите в Таблицу 2 файла протокола.
\end{enumerate}